\documentclass[11pt]{article}

\usepackage[margin=1in]{geometry}
\usepackage[utf8]{inputenc}
\usepackage[T1]{fontenc}
\usepackage{lmodern}
\usepackage{amsmath,amssymb,amsthm}
\usepackage{booktabs,array}
\usepackage{microtype}
\usepackage[numbers,sort&compress]{natbib}
\usepackage{xcolor}
\usepackage{etoolbox}
\AtBeginEnvironment{table}{\small\setlength{\tabcolsep}{3pt}}
\usepackage[colorlinks=true,linkcolor=blue!50!black,citecolor=blue!50!black,urlcolor=blue!50!black]{hyperref}

\newcommand{\draftnote}[1]{\textcolor{red!70!black}{[\textbf{Draft:} #1]}}

\title{Learning Finite-State Transducers Exactly by Gradient Descent, with Proof\\[0.5em]
\large Skeleton, 25 October 2026: not for circulation}

\author{
  David Barkhuizen\\
  Independent researcher
  \and
  Claude (Opus 5.5)\\
  Anthropic
}

\date{}

\begin{document}

\maketitle

\draftnote{The paper workplan's stage 6 (indrajala-ml \texttt{docs/workplans/paper-workplan.md}, D1
and D7): built from \texttt{snapped-window-recurrence.tex} (P1) and
\texttt{certified-modular-division.tex} (P2), which stay until this supersedes them. Every number in a
table is regenerated by indrajala-ml's \texttt{scripts/paper\_tables.py} (T7); \texttt{paper\_tables.py
check} compares this file's tables with the data. Waiting on: T10 (D19) and T11 (D20), whose
batches are running; K7 is written after T11's reading. The authorship and contribution statement
follow D6.}

\begin{abstract}
\draftnote{Written last. Its points: networks trained by gradient descent on input--output examples
alone (never states or traces) learn finite-state transducers of up to 169 states, each proved equal
to its target at every length by a complete certificate; sampled tests overstate this; the snap
makes the learned machine extrapolate; training on the target's approximations makes product
machines learnable (16 of 20 against 0 of 20 at 121 states), and because they are approximations;
the learned factors, read exactly, hold the target's approximations; and the proof needs no target
in the loop.}
\end{abstract}

\section{Introduction}
\label{sec:intro}

\draftnote{The question (D1): can a network trained by gradient descent on input--output examples
alone come to compute a finite-state transducer exactly at every length, how large, and how would
one know? The gap: extracted-machine work stops at a few states and checks by sampling or queries
\citep{giles-1992-learning-extracting-fsa-second-order-rnn,zeng-1993-self-clustering-recurrent-networks,weiss-2018-extracting-automata-queries-counterexamples};
classical learners find these machines from far fewer examples
\citep{oncina-1992-transducer-learning-pattern-recognition,angluin-1987-learning-regular-sets-queries-counterexamples},
so the question is about neural sequence models trained as they are trained. Contributions in the
claims' order, K1--K9 (Section~\ref{sec:results}), each with its nearest prior work.}

\section{Related work}
\label{sec:related}

\draftnote{From P1's and P2's related work and T8's sweep (indrajala-research's
\texttt{window-recurrent-literature.md}, its 2026-10-23 extension: no 2025--2026 work learns an
automaton by gradient descent and proves it at every length). Paragraphs: automaton extraction and
discretised states; verification of learned automata; state tracking and length generalisation
\citep{deletang-2022-neural-networks-chomsky-hierarchy,merrill-2024-illusion-of-state-ssm};
auxiliary outputs and intermediate supervision (P2's); algebraic decompositions
\citep{krohn-rhodes-1965-algebraic-theory-machines,egri-nagy-2024-covering-lemma-transformation-semigroups};
conformance testing \citep{chow-1978-testing-software-design-finite-state-machines}.}

\section{Method}
\label{sec:method}

\paragraph{Targets.} A target is a Mealy machine read one column at a time, its output asked at every
column. \emph{Division} at $p$ reads two base-3 numbers $A$ and $B$, most significant digit first, a
column's token $3a_i + b_i$, and outputs the running $(A \cdot B^{-1}) \bmod p$, or ``undefined''
while $B \equiv 0$; its minimal machine has $p^2$ states (25, 49, 121, 169 at $p = 5, 7, 11, 13$),
the product of two modular counters. The \emph{algebras}: the word problems of $S_3$ and $S_4$ (6 and
24 states), move-to-front on 4 items (24 states), and random machines of 12 and 16 states.

\paragraph{Approximations.} An \emph{approximation} of a target is a homomorphic image of it, a
quotient by a congruence of its states, as
\citet{egri-nagy-2024-covering-lemma-transformation-semigroups} use the term. Division's include
the running sum $(A + B) \bmod p$ and difference $(A - B) \bmod p$; for odd $p$ the two together
determine the state, so their product \emph{covers} the target and they are its
\emph{coordinates}. The quotient and the product $AB \bmod p$ are no approximations. A network trained
on approximations answers them on the same input through query tokens, one output layer for all.

\paragraph{The network.} \draftnote{P2's network section, at the paper's cell (D12): a window of
columns read by a causal transformer, the state handed on by gated recurrent cells whose outputs are
snapped to the nearest of $m$ learned codes, straight-through, in two factors (each its own pooling,
cell and codebook, every block reading both); the factors' codes, $\lceil 32\sqrt{q}/11 \rceil$
each, 32 wide; reading the output; the fork and the commitment as built; the 32-wide window. T1 and
T3 use the approximation claim's cell (two factors of 16, the default window, 400 epochs, no fork).}

\paragraph{The certificate.} \draftnote{P1's: every window from every reachable joint code mapped onto
the target, every output and next state checked, the learned machine minimised; \emph{certified}
with the network's own snaps, \emph{robustly certified} with every near tie (snap ratio below 1.05)
branched both ways. Every network counted here passes both. Long tests: uniform and stratified
catalogues at 128, 256 and 1{,}024 columns, 375{,}000 cases a network.}

\paragraph{Training and feasibility.} \draftnote{The curriculum (step $j$ trains on lengths up to
$j$, warm-started; a step ends at its check, a plateau or its epochs), stopped at the first
certified step; the epoch budget 600 (400 for T1 and T3). Every cell's codes at least 3 times the
target's states, its longest training length the exclusion bound $m + q - 1$; tested against the
config.}

\paragraph{Statistics.} Every comparison at 20 fresh seeds an arm, registered with its direction
before its data; Fisher's exact test, one-sided as registered; Holm across the paper's family;
Wilson 95\% intervals. A run counts as a success only when certified in both verdicts and every
long-test case passed.

\section{Results}
\label{sec:results}

\subsection{K1: a snapped network is a machine, and a certificate proves it}
\draftnote{From P1: the certificate on 724 networks; every certified network passing every long-test
case. Here: the paper's runs, 60 of T2 and T5 alone passing 34{,}614{,}000 long-test cases.}

\subsection{K2: sampled tests overstate}
\draftnote{From P1's E1 and E3 (registered, Holm): networks passing every sampled test, one passing
328{,}000 long-test cases, provably wrong; more codes buy test-passing networks, not correct ones.
The phenomenon is \citet{weiss-2018-extracting-automata-queries-counterexamples}'s; the complete
check is ours.}

\subsection{K3: the snap is what extrapolates}
\label{sec:baselines}
Division at $p = 5$ trained on its cover, three arms that differ in one thing each from the first
(T3, registered):

\begin{table}[t]
\centering
\caption{The snap and the transformer at $p = 5$ (T3): successes of 20; failures that came near a solution (every answer at least 0.95 at some diagnosis); failures that collapsed.}
\label{tab:baselines}
\begin{tabular}{lcccc}
\toprule
arm & succeeded & Wilson 95\% & near & collapsed \\
\midrule
snapped codes, transformer window (control) & 13 of 20 & 0.43--0.82 & 6 of 7 & 5 \\
continuous state (no codes) & 0 of 20 & 0.00--0.16 & 20 of 20 & 0 \\
linear window (no blocks) & 0 of 20 & 0.00--0.16 & 0 of 20 & 8 \\
\bottomrule
\end{tabular}
\end{table}

The control certified more often than the continuous state (one-sided $p = 6.44 \times 10^{-6}$,
Holm $1.29 \times 10^{-5}$) and than the linear window (the same). Every continuous network came
near a solution on the training lengths and failed every long-test case from 128 columns: a
continuous state fits the training lengths and drifts beyond them.

\subsection{K4 and K6: approximations make product machines learnable, at scale}
\label{sec:scale}

\begin{table}[t]
\centering
\caption{Division at scale, the paper's cell (T2): successes of 20, the successes' epochs to certify
(range, median), and failures that came near a solution.}
\label{tab:scale}
\begin{tabular}{lccccc}
\toprule
backdrop & $q$ & succeeded & Wilson 95\% & epochs & near \\
\midrule
\texttt{d7-family} & 49 & 17 of 20 & 0.64--0.95 & 81--474, 195 & 3 of 3 \\
\texttt{d11-family} & 121 & 16 of 20 & 0.58--0.92 & 146--430, 261 & 2 of 4 \\
\texttt{d13-family} & 169 & 15 of 20 & 0.53--0.89 & 187--534, 313 & 4 of 5 \\
\texttt{d11-quotient} & 121 & 0 of 20 & 0.00--0.16 & -- & 0 of 20 \\
\bottomrule
\end{tabular}
\end{table}

\draftnote{K6: the registered bar (a Wilson lower bound of 0.25) met at every $p$. K4 at 121 states:
the cover 16 of 20 against the quotient alone 0 of 20, one-sided $p = 7.7 \times 10^{-8}$; P2's
registered $p = 5$ comparison (13 of 20 against 1 of 20) beside it.}

\subsection{K5: one approximation suffices, because it is one}
\label{sec:approximation}

\begin{table}[t]
\centering
\caption{The approximation claim at $p = 5$ (P2): certified of 20, the successes' steps and epochs.}
\label{tab:approximation}
\begin{tabular}{lccc}
\toprule
extra outputs & certified & Wilson 95\% & steps, epochs \\
\midrule
sum, difference (cover) & 15 of 20 & 0.53--0.89 & 6--12, 78--235 \\
sum (an approximation) & 10 of 20 & 0.30--0.70 & 7--11, 64--160 \\
none & 1 of 20 & 0.01--0.24 & 12, 279 \\
product (not an approximation) & 0 of 20 & 0.00--0.16 & -- \\
\bottomrule
\end{tabular}
\end{table}

\begin{table}[t]
\centering
\caption{The approximation against two learnable extra outputs at $p = 5$ (T1): successes of 20; failures near a solution; failures that collapsed.}
\label{tab:extra-outputs}
\begin{tabular}{lcccc}
\toprule
extra output & succeeded & Wilson 95\% & near & collapsed \\
\midrule
sum (an approximation) & 6 of 20 & 0.15--0.52 & 11 of 14 & 10 \\
suffix (state-carrying, no approximation) & 0 of 20 & 0.00--0.16 & 0 of 20 & 2 \\
token (stateless) & 1 of 20 & 0.01--0.24 & 0 of 19 & 7 \\
\bottomrule
\end{tabular}
\end{table}

\draftnote{T1a: the sum more often than the suffix, one-sided $p = 0.0101$ (Holm 0.0202); T1b: than
the token, 0.0457 (Holm 0.0457). An extra output the network can learn does not help as the
approximation does; the state-carrying one never certified. The sum's 6 of 20 below A2's 10 of 20,
stated.}

\subsection{K7: beyond arithmetic}
\label{sec:breadth}

\begin{table}[t]
\centering
\caption{The algebras and random machines, the paper's cell (T5): successes of 20.}
\label{tab:breadth}
\begin{tabular}{lccccc}
\toprule
backdrop & $q$ & succeeded & Wilson 95\% & epochs & near \\
\midrule
\texttt{s3} & 6 & 1 of 20 & 0.01--0.24 & 52 & 0 of 19 \\
\texttt{s4} & 24 & 4 of 20 & 0.08--0.42 & 157--394, 191 & 6 of 16 \\
\texttt{mtf4} & 24 & 4 of 20 & 0.08--0.42 & 152--309, 195.5 & 1 of 16 \\
\texttt{rq12} & 12 & 3 of 20 & 0.05--0.36 & 95--230, 120 & 8 of 17 \\
\texttt{rq16} & 16 & 0 of 20 & 0.00--0.16 & -- & 11 of 20 \\
\bottomrule
\end{tabular}
\end{table}

\draftnote{K7 as existence at low rates (T5a): $S_3$ refuted at the bar, $S_4$ and move-to-front
undecided; random machines certify at 12 states, none at 16. Waits on T11 (the algebras given their
approximations) and T10 (a planted approximation against primitive machines).}

\subsection{K8: the learned factors hold the target's approximations}
\label{sec:decomposition}

\begin{table}[t]
\centering
\caption{The decomposition reading (T4, registered before the reading): networks; factors whose
approximation is nontrivial; networks naming a congruence in a nontrivial factor (in both);
covers; congruences determined.}
\label{tab:decomposition}
\begin{tabular}{lcccc>{\raggedright\arraybackslash}p{2.9cm}}
\toprule
set & networks & nontrivial & naming (both) & covers & determines \\
\midrule
$p$ 5 division claim, family & 13 & 8/26 & 5 (3) & 3 & \texttt{A + B} 4, \texttt{A - B} 4 \\
$p$ 5 division claim, quotient & 1 & 0/2 & 0 (0) & 0 & -- \\
$p$ 5 approximation claim, family & 15 & 13/30 & 8 (5) & 5 & \texttt{A + B} 7, \texttt{A - B} 6 \\
$p$ 5 approximation claim, sum & 10 & 7/20 & 7 (0) & 0 & \texttt{A + B} 7 \\
$p$ 5 approximation claim, quotient & 1 & 1/2 & 1 (0) & 0 & \texttt{B} 1 \\
$p$ 5 probe 2a (apart) & 3 & 3/6 & 2 (1) & 1 & \texttt{A + B} 2, \texttt{A - B} 1 \\
$p$ 7 probe 4 and fork probe & 6 & 3/12 & 2 (1) & 1 & \texttt{A - B} 2, \texttt{A + B} 1 \\
$p$ 11 ladder & 4 & 6/8 & 4 (2) & 2 & \texttt{A + B} 3, \texttt{A - B} 3 \\
$p$ 13 W1 & 5 & 2/10 & 2 (0) & 0 & \texttt{A + B} 1, \texttt{A - B} 1 \\
word-problem-s3 & 1 & 2/2 & 1 (1) & 1 & \texttt{places \{1\}} 1, \texttt{places \{2\}} 1 \\
word-problem-s4 & 4 & 3/8 & 3 (0) & 0 & \texttt{places \{1\}} 2, \texttt{places \{2\}} 1 \\
move-to-front-3 & 6 & 9/12 & 6 (3) & 1 & \texttt{order \{0, 2\}} 5, \texttt{order \{0, 1\}} 4, \texttt{order \{1, 2\}} 1 \\
\bottomrule
\end{tabular}
\end{table}

\begin{table}[t]
\centering
\caption{The decomposition reading on T2's and T5's fresh networks.}
\label{tab:decomposition-fresh}
\begin{tabular}{lcccc>{\raggedright\arraybackslash}p{2.9cm}}
\toprule
backdrop & networks & nontrivial & naming (both) & covers & determines \\
\midrule
\texttt{d7-family} & 17 & 9/34 & 8 (1) & 1 & \texttt{A - B} 7, \texttt{A + B} 2 \\
\texttt{d11-family} & 16 & 9/32 & 8 (1) & 1 & \texttt{A - B} 6, \texttt{A + B} 3 \\
\texttt{d13-family} & 15 & 6/30 & 5 (1) & 1 & \texttt{A - B} 4, \texttt{A + B} 2 \\
\texttt{s3} & 1 & 0/2 & 0 (0) & 0 & -- \\
\texttt{s4} & 4 & 1/8 & 1 (0) & 0 & \texttt{places \{3\}} 1 \\
\texttt{mtf4} & 4 & 8/8 & 4 (1) & 0 & \texttt{order \{1, 3\}} 3, \texttt{order \{0, 3\}} 2, \texttt{order \{0, 1, 3\}} 1, \texttt{order \{0, 1\}} 1, \texttt{order \{1, 2\}} 1 \\
\texttt{rq12} & 3 & 0/6 & 0 (0) & 0 & -- \\
\bottomrule
\end{tabular}
\end{table}

\draftnote{T4a and T4b held on both sets (the algebras' factors 14 of 22, then 9 of 18, nontrivial;
no factor holding the sign, every congruence named on $S_3$ and $S_4$ a stabiliser's); T4c mixed on
both (division 43 of 116, then 24 of 96): a minority of division's factors decompose, each into
exactly one line, always a trained one.}

\subsection{K9: the proof needs no target in the loop}
\label{sec:proof}

\begin{table}[t]
\centering
\caption{The answers-only stop and Chow's test set on every coded run of T1, T2, T3 and T5 (T6): 280
runs at the exact state bound and at 1.1 times it.}
\label{tab:proof}
\begin{tabular}{lcc}
\toprule
outcome & exact bound & $1.1 \times$ bound \\
\midrule
no stop (the run never certified) & 200 & 200 \\
kept, Chow proved & 80 & 43 \\
kept, out of Chow's reach & 0 & 17 \\
false stop, Chow proved the machine equal to the target & 0 & 5 \\
false stop, Chow refuted & 0 & 4 \\
false stop, out of Chow's reach & 0 & 11 \\
\bottomrule
\end{tabular}
\end{table}

\begin{table}[t]
\centering
\caption{Stops at $1.1 \times$ the bound that Chow's test set refuted: the bound, the stop's and the
certified step, the counterexample's length in tokens.}
\label{tab:refuted}
\begin{tabular}{lccc}
\toprule
run & bound & stop, certified & counterexample \\
\midrule
\texttt{d7-family-300} & 53 & 7, 11 & 5 \\
\texttt{d11-family-300} & 133 & 17, 22 & 7 \\
\texttt{d11-family-310} & 133 & 13, 32 & 5 \\
\texttt{d13-family-316} & 185 & 10, 11 & 4 \\
\bottomrule
\end{tabular}
\end{table}

\draftnote{T6a and T6c held: at the exact bound no false stop in 280 runs, every certified run's
stop at its certified step, and Chow's test set \citep{chow-1978-testing-software-design-finite-state-machines}
proved each of the 80. T6b refuted as registered: at 1.1 the stop rule landed on 4 wrong machines,
each refuted by Chow; nothing false was proved. A discoverer who knows the size only roughly needs
the proof in the loop.}

\subsection{K10: training past a proof can destroy it}
\draftnote{Short, stated as exploratory: P1's replays and P2's collapses.}

\section{Discussion}
\draftnote{Limits: arithmetic carries the scale; the algebras at low rates; primitive random
machines past 12 states; the certificate in the loop (K9 answers it offline).}

\section{Reproducibility}
\draftnote{T7: every table regenerated from the data by \texttt{paper\_tables.py}; the bit-for-bit
reruns (D21); the datapack's manifests; the repositories public (D6).}

\section*{Author contributions}
\draftnote{D6: the owner's direction and decisions; Claude's design, code, runs, analysis and
drafting.}

\bibliographystyle{plainnat}
\bibliography{library,extra}

\end{document}
